% !TEX root =./ECE448_Practice_main.tex
%
% L2 practice set (Noise Figure and Link Budget). Blank copy: \begin{solution}[h]
% reserves h of work space. SOLUTIONS copy (\iskey=1) prints the red worked solutions.

\fullwidth{\textbf{Learning Objective 1.1: } I can define noise figure and explain
how each stage of a receiver chain, the first above all, sets the system's noise performance.
\\[2pt]\textbf{Learning Objective 1.2: } I can compute a link budget for a specified
geometry, transmitter, and antenna configuration, and compare it with the signal an SDR shows.
\\[4pt]\normalfont\small Show your work. Write a \textbf{Documentation} line at the
end (\textbf{None} if you did not collaborate).}

\begin{questions}

%%%%%%%%%%%%%%%%%%%%%%%%%%%%%%%%%%%%%%%%%
\question \textbf{Decibels and the ISM bands.} A small 915 MHz telemetry radio
puts out 100 mW. It connects to a 9 dBi antenna through a cable and connectors
with 2 dB of loss.
\begin{parts}

	\begin{minipage}[t]{0.58\linewidth}
		\part Express 5 W and 0.1 $\mu$W in dBm.
		\begin{solution}[1.0in]
		\eq{ 10\mylog{\dfrac{5000\ \text{mW}}{1\ \text{mW}}} &= 37.0\ \text{dBm} \\
			10\mylog{\dfrac{10^{-4}\ \text{mW}}{1\ \text{mW}}} &= -40.0\ \text{dBm} }
		\end{solution}
	\end{minipage}\hfill%
	\begin{minipage}[t]{0.37\linewidth}
		\raggedleft \vspace{12pt}
		5 W = \ansbox[1.3in]{$37.0$ dBm}

		\vspace{4pt}
		0.1 $\mu$W = \ansbox[1.3in]{$-40.0$ dBm}
	\end{minipage}

	\begin{minipage}[t]{0.58\linewidth}
		\part Find the radio's EIRP in dBm and in milliwatts.
		\begin{solution}[1.0in]
		\eq{ \text{EIRP} &= 10\mylog{100} - 2 + 9 = 20 - 2 + 9 = 27\ \text{dBm} \\
			&= 10^{27/10} = 501\ \text{mW} }
		\end{solution}
	\end{minipage}\hfill%
	\begin{minipage}[t]{0.37\linewidth}
		\raggedleft \vspace{12pt}
		EIRP = \ansbox[1.3in]{$27$ dBm}

		\vspace{4pt}
		EIRP = \ansbox[1.3in]{$501$ mW}
	\end{minipage}

	\part Which of the course SDRs (RTL-SDR, PlutoSDR, HackRF One) can receive the
	2.4 GHz ISM band, and which can transmit in it?
		\begin{solution}[0.9in]
		The PlutoSDR (325 MHz to 3.8 GHz) and the HackRF One (1 MHz to 6 GHz) both
		cover 2.4 GHz and can both transmit. The RTL-SDR tops out near 1.77 GHz and is
		receive-only, so it can reach the 915 MHz band but not 2.4 GHz.
		\end{solution}

	\part Explain in one or two sentences why Wi-Fi, Bluetooth, and most hobby
	drones all operate in the same few bands.
		\begin{solution}[0.9in]
		The ISM bands are the bands where FCC Part 15 allows low-power transmitters to
		operate without a license, as long as they accept interference from other
		users. Every consumer device that cannot require its user to hold a license
		ends up there.
		\end{solution}

\end{parts}

\newpage
%%%%%%%%%%%%%%%%%%%%%%%%%%%%%%%%%%%%%%%%%
\question \textbf{Free-space path loss.} A 915 MHz radio talks to a
ground station 10 km away. Both antennas have 6 dBi of gain, and the transmitter
delivers 1 W to its antenna.
\begin{parts}

	\begin{minipage}[t]{0.58\linewidth}
		\part Find the free-space path loss at 915 MHz and 10 km.
		\begin{solution}[1.0in]
		\eq{ L_p &= 20\mylog{915\times10^{6}} + 20\mylog{10^{4}} - 147.55 \\
			&= 179.2 + 80.0 - 147.55 = 111.7\db }
		\end{solution}
	\end{minipage}\hfill%
	\begin{minipage}[t]{0.37\linewidth}
		\raggedleft \vspace{12pt}
		$L_p$ = \ansbox[1.3in]{$111.7\db$}
	\end{minipage}

	\begin{minipage}[t]{0.58\linewidth}
		\part Find the received power in dBm.
		\begin{solution}[1.0in]
		\eq{ P_r &= P_t + G_t + G_r - L_p \\
			&= 30 + 6 + 6 - 111.7 = -69.7\ \text{dBm} }
		\end{solution}
	\end{minipage}\hfill%
	\begin{minipage}[t]{0.37\linewidth}
		\raggedleft \vspace{12pt}
		$P_r$ = \ansbox[1.3in]{$-69.7$ dBm}
	\end{minipage}

	\begin{minipage}[t]{0.58\linewidth}
		\part The same system moves to 2.45 GHz with the same antenna gains and
		twice the distance. Find the new path loss using the rules of thumb, then
		check it with the full expression.
		\begin{solution}[1.3in]
		Doubling the distance adds 6.0 dB. The frequency ratio adds
		$20\mylog{2450/915} = 8.6$ dB. So
		\eq{ L_p &\approx 111.7 + 6.0 + 8.6 = 126.3\db \\
			L_p &= 20\mylog{2.45\times10^{9}} + 20\mylog{2\times10^{4}} \\
			&\quad - 147.55 = 126.3\db }
		\end{solution}
	\end{minipage}\hfill%
	\begin{minipage}[t]{0.37\linewidth}
		\raggedleft \vspace{12pt}
		$L_p$ = \ansbox[1.3in]{$126.3\db$}
	\end{minipage}

	\part Path loss is not absorption, yet it grows with frequency. Explain why in
	one or two sentences.
		\begin{solution}[0.9in]
		The power spreads over a sphere regardless of frequency, but an antenna of
		fixed gain has an effective area proportional to $\lambda^2$. At a higher
		frequency the same gain collects power over a smaller area, so less of the
		spreading wave reaches the receiver.
		\end{solution}

\end{parts}

\newpage
%%%%%%%%%%%%%%%%%%%%%%%%%%%%%%%%%%%%%%%%%
\question \textbf{The noise floor.} An RTL-SDR with a 6 dB noise figure receives a
broadcast FM station. One FM channel occupies about 200 kHz.
\begin{parts}

	\begin{minipage}[t]{0.58\linewidth}
		\part Find the receiver's noise floor in one 200 kHz channel.
		\begin{solution}[1.0in]
		\eq{ N &= -174 + 10\mylog{200\times10^{3}} + 6 \\
			&= -174 + 53.0 + 6 = -115.0\ \text{dBm} }
		\end{solution}
	\end{minipage}\hfill%
	\begin{minipage}[t]{0.37\linewidth}
		\raggedleft \vspace{12pt}
		$N$ = \ansbox[1.3in]{$-115.0$ dBm}
	\end{minipage}

	\begin{minipage}[t]{0.58\linewidth}
		\part The station arrives at $-80$ dBm. Find the SNR in the channel.
		\begin{solution}[0.8in]
		\eq{ \text{SNR} = -80 - (-115.0) = 35.0\db }
		\end{solution}
	\end{minipage}\hfill%
	\begin{minipage}[t]{0.37\linewidth}
		\raggedleft \vspace{12pt}
		SNR = \ansbox[1.3in]{$35.0\db$}
	\end{minipage}

	\part A student leaves the receiver filter open to the full 2.4 MHz sample
	rate and measures the station's power against all of that noise. Why does the
	wide filter show a lower SNR, and which SNR does the FM demodulator experience?
		\begin{solution}[1.1in]
		Noise power grows with bandwidth, but the station's power does not: it all lies
		inside its 200 kHz channel. The wide filter lets in about 12 times as much noise
		and no more signal, so the SNR looks lower. A demodulator that filters to the
		channel first sees the 35.0 dB of part (b).
		\end{solution}

\end{parts}

\newpage
%%%%%%%%%%%%%%%%%%%%%%%%%%%%%%%%%%%%%%%%%
\question \textbf{Where the LNA goes.} A ground station has a 2 dB cable run, an
LNA with NF 0.8 dB and 18 dB of gain, and an SDR with NF 7 dB. The system noise
figure of each arrangement is:

\begin{center}
\begin{tabular}{ll}
	antenna $\to$ cable $\to$ LNA $\to$ SDR & 3.0 dB \\
	antenna $\to$ LNA $\to$ cable $\to$ SDR & 1.2 dB \\
	antenna $\to$ cable $\to$ SDR (no LNA) & 9.0 dB \\
\end{tabular}
\end{center}

\begin{parts}

	\part Where should the LNA go, and what would you tell someone who wants to mount
	it next to the SDR to keep it out of the weather?
		\begin{solution}[1.0in]
		At the antenna, ahead of the cable: that arrangement has the lowest noise
		figure, 1.2 dB. Mounted next to the SDR it gives 3.0 dB, so the receiver's
		noise floor sits 1.8 dB higher. Use a weatherproof LNA at the antenna instead.
		\end{solution}

	\part Explain why the cable's loss matters so much more in front of the LNA than
	behind it.
		\begin{solution}[1.0in]
		In front of the LNA the cable is the first stage, and nothing ahead of it
		divides its noise down, so its 2 dB adds almost directly to the noise figure.
		Behind the LNA, its 18 dB of gain sits ahead of both the cable and the SDR, so
		the cable changes the total by only about 0.2 dB.
		\end{solution}

	\part Without the LNA the noise figure is 9.0 dB, the cable's 2 dB plus the SDR's
	7 dB. You add the LNA at the antenna. What happens on the SDR's display to the
	noise floor, to a weak signal near it, and to the SNR?
		\begin{solution}[1.1in]
		Both rise on the display: the signal by the LNA's 18 dB, and the floor by only
		about 10 dB, because the system adds 7.8 dB less noise of its own. The weak
		signal stands about 8 dB clearer of the floor, so the SNR improves by about
		8 dB.
		\end{solution}

\end{parts}

\newpage
%%%%%%%%%%%%%%%%%%%%%%%%%%%%%%%%%%%%%%%%%
\question \textbf{A GPS link budget.} A GPS satellite transmits the L1 signal at
1575.42 MHz with an EIRP of 27 dBW. A receiver on the ground sees the satellite
directly overhead at 20,200 km through a 0 dBi antenna. Polarization and
atmospheric losses total 2 dB. The receiver has a 2 dB noise figure and a 2 MHz
bandwidth.
\begin{parts}

	\begin{minipage}[t]{0.58\linewidth}
		\part Find the free-space path loss and the received power in dBm.
		\begin{solution}[1.5in]
		\eq{ L_p &= 20\mylog{1.57542\times10^{9}} \\
			&\quad + 20\mylog{2.02\times10^{7}} - 147.55 \\
			&= 183.9 + 146.1 - 147.55 = 182.5\db \\
			\text{EIRP} &= 27\ \text{dBW} = 57\ \text{dBm} \\
			P_r &= 57 - 182.5 + 0 - 2 = -127.5\ \text{dBm} }
		\end{solution}
	\end{minipage}\hfill%
	\begin{minipage}[t]{0.37\linewidth}
		\raggedleft \vspace{12pt}
		$L_p$ = \ansbox[1.3in]{$182.5\db$}

		\vspace{4pt}
		$P_r$ = \ansbox[1.3in]{$-127.5$ dBm}
	\end{minipage}

	\begin{minipage}[t]{0.58\linewidth}
		\part Find the noise floor and the SNR at the receiver.
		\begin{solution}[1.1in]
		\eq{ N &= -174 + 10\mylog{2\times10^{6}} + 2 = -109.0\ \text{dBm} \\
			\text{SNR} &= -127.5 - (-109.0) = -18.5\db }
		\end{solution}
	\end{minipage}\hfill%
	\begin{minipage}[t]{0.37\linewidth}
		\raggedleft \vspace{12pt}
		$N$ = \ansbox[1.3in]{$-109.0$ dBm}

		\vspace{4pt}
		SNR = \ansbox[1.3in]{$-18.5\db$}
	\end{minipage}

	\part The signal arrives below the noise floor, yet GPS works everywhere.
	Explain in one or two sentences how.
		\begin{solution}[0.9in]
		The receiver already knows each satellite's code, so it correlates the received
		signal against that code. The code adds up coherently over its whole length while the
		noise does not, which pulls the signal back out of the noise.
		\end{solution}

	\part A car's GPS antenna sits on the roof with 5 m of cable to the dashboard
	receiver. Where should the LNA go, and why? Use your answer to (b) to explain
	what is at stake.
		\begin{solution}[1.1in]
		The LNA belongs at the antenna, ahead of the cable; that is why GPS antennas with
		a built-in LNA, active antennas, are common. The signal already arrives 18.5 dB
		below the noise floor, so every dB of cable loss in front of the first amplifier
		raises the floor and comes straight out of what the correlation has to work with.
		\end{solution}

\end{parts}

\newpage
%%%%%%%%%%%%%%%%%%%%%%%%%%%%%%%%%%%%%%%%%
\question \textbf{Reading SNR on an SDR.} A GNU Radio frequency sink shows an
unmodulated carrier from a signal generator, received by an RTL-SDR running at
$f_s = 2.4$ MHz with an FFT size of 2048. The carrier's peak reads $-32$ dBFS and
the noise floor beside it reads $-78$ dBFS.
\begin{parts}

	\begin{minipage}[t]{0.58\linewidth}
		\part Find the SNR the display shows, and the width of one FFT bin.
		\begin{solution}[1.0in]
		\eq{ \text{SNR} &= -32 - (-78) = 46\db \\
			\dfrac{f_s}{N} &= \dfrac{2.4\ \text{MHz}}{2048} = 1.17\ \text{kHz} }
		\end{solution}
	\end{minipage}\hfill%
	\begin{minipage}[t]{0.37\linewidth}
		\raggedleft \vspace{12pt}
		SNR = \ansbox[1.3in]{$46\db$}

		\vspace{4pt}
		bin = \ansbox[1.3in]{$1.17$ kHz}
	\end{minipage}

	\part The display reads dBFS, not dBm. Why can you still read an SNR off it?
		\begin{solution}[0.8in]
		SNR is a difference between two levels on the same axis, so the unknown offset
		between dBFS and dBm cancels.
		\end{solution}

	\part You change the FFT size from 2048 to 8192. What happens to the noise floor,
	to the carrier's peak, and to the SNR the display shows, and why?
		\begin{solution}[1.1in]
		Each bin gets narrower and holds less of the noise, so the floor drops. The
		carrier's power is all in one bin, so its peak stays where it was, and the
		displayed SNR rises, though the signal itself did not change.
		\end{solution}

\end{parts}

\vspace{1.5em}
\noindent\textbf{Documentation:} \rule[-2pt]{0.6\linewidth}{0.4pt}

\end{questions}
